\section{Introduction}\label{sec:intro}

Kinetic reduced magnetohydrodynamics (KRMHD), the limit $k_\perp\rho_i\to0$ of gyrokinetics derived by \citet{Schekochihin2009}, hereafter S09, describes anisotropic fluctuations at scales much larger than the ion gyroradius $\rho_i$ in a strongly magnetised, weakly collisional plasma. Its Alfv\'enic fluctuations obey reduced magnetohydrodynamics and advect the compressive fluctuations, the slow waves and the entropy mode, which obey a kinetic equation for the ion distribution in the parallel velocity $v_\parallel$ and do not act back on them. Without collisions the compressive fluctuations are damped \citetext{\citealp{Barnes1966}; S09} by phase mixing, which transfers their free energy to ever finer structure in $v_\parallel$; this transfer is central to turbulent heating \citep[S09;][]{Schekochihin2016,Meyrand2019,Zhou2023,Adkins2025}. In a Hermite expansion in $v_\parallel$, phase mixing is a flux of free energy to large Hermite order $m$ \citep{Kanekar2015,Schekochihin2016,AdkinsSchekochihin2018,Meyrand2019,Kawazura2020}, and linear phase mixing driven at low $m$ gives a spectrum proportional to $m^{-1/2}$ \citep{ZoccoSchekochihin2011,Kanekar2015}. In KRMHD turbulence \citet{Meyrand2019} found a return flux that cancels much of the phase-mixing flux on average and attributed it to stochastic plasma echoes, in which nonlinear advection returns free energy to low $m$ \citep{KanekarThesis2015,Schekochihin2016,Parker2016,AdkinsSchekochihin2018}.

\citet{Chandran2026}, hereafter CMM, showed that gyrokinetics conserves a second three-dimensional quadratic invariant, the gyrokinetic helicity, and derived its form in KRMHD, where the compressive sector splits into two kinetic branches $G^\pm$ and the Laguerre components of order $l\ge2$ in $v_\perp$ (S09, \S6.2; CMM, appendix~B). In each branch the parallel electric field and the mirror force couple the distribution to its zeroth velocity moment through a term of strength $1/\Lambda^\pm$, which depends only on $\beta_i$ and $T_i/(ZT_e)$; we call this term the Landau coupling. Each branch conserves its free energy $W^\pm$ (S09) and, as CMM found, the phase-space part $\Hph^\pm$ of the helicity and a further quadratic invariant $\Gamma^\pm$. CMM asked whether helicity conservation leads to a bottleneck in velocity space.

We show that the compressive sector of KRMHD is exactly a continuum of independent passive scalars labelled by phase velocity. Because the field relations of KRMHD are algebraic, the branch operator, which combines streaming with the Landau coupling, contains no wavenumber. It is self-adjoint in the free-energy inner product, and its Case--Van Kampen transform \citep{vanKampen1955,Case1959} acts on velocity alone, so it commutes with both nonlinear terms. The slice $a_u$, the transform at $u=v_\parallel/\vth$ with $\vth$ the ion thermal speed, obeys
\begin{equation*}
  \partial_ta_u+\big(\uE+u\vth\,\bhat\big)\cdot\nabla a_u=0,
\end{equation*}
where $\uE$ is the $\boldsymbol E\times\boldsymbol B$ velocity and $\bhat$ the unit vector along the perturbed field. Slices with different $u$ do not interact, in three dimensions and for any smooth Alfv\'enic fields; only velocity-dependent collisions couple them (section~\ref{sec:slices}). The energy of every slice is conserved. The free energy and CMM's $\Gamma$ and $\Hph$ are the integrals of the slice energies over $u$ weighted by $1$, $u$ and $1/u$, any other weight gives a further invariant, and $\Gamma$ is the net free-energy flux along the guide field (section~\ref{sec:invariants}). Phase mixing is the decorrelation of slices by their relative streaming along the field, and slices decorrelated in velocity give the $m^{-1/2}$ spectrum of forced linear phase mixing. Linear phase mixing conserves every invariant of the compressive distribution while it moves all of the free energy to arbitrarily large $m$, so neither the helicity nor any other invariant imposes a constraint of the kind \citet{Fjortoft1953} found for two-dimensional flows, as CMM expected for $\Gamma$ (section~\ref{sec:decorrelation}); this concerns transient solutions, not the forced steady states in which CMM posed their question. A Hermite hierarchy truncated at order $M$, as evolved by Fourier--Hermite codes, consists of exactly $M+1$ slices, which for the closure $g_{M+1}=0$ sit at the Gauss nodes of the Case--Van Kampen spectral measure; the fastest slice sets the time step of an explicit scheme, and the phase-space helicity of the retained moments is conserved only at odd $M$ (section~\ref{sec:truncated}). Section~\ref{sec:numerics} checks these results with the GANDALF code \citep{Kanekar2025}.

Parts of this structure are known elsewhere. The transform is, up to a weight, the inverse of the G-transform of \citet{MorrisonPfirsch1992} and \citet{Morrison1994}, which \citet{Morrison2000} showed to be invertible under the stability condition of \citet{Penrose1960}, and which \citet{SkiffHowes2024} use to infer the electric field over a region from the velocity distribution measured at one point. \citet{Heninger2018} applied it to electrostatic drift kinetics with Boltzmann electrons, where the transformed components have the structure of our slices but remain coupled through the potential (their eq.~85); the energy of each component is conserved there too, although they do not state this. Invariants at each velocity were noted for a linear problem by \citet{PlunkParker2014} and proved for a quasi-two-dimensional nonlinear model by \citet{Plunk2015}, and \citet[\S5.1]{Kunz2015} remarked that the free energy of their pressure-anisotropic generalisation of KRMHD is one of infinitely many invariants, without specifying the others. For the Laguerre components with $l\ge2$, which have no Landau coupling, S09's eqs~(210) and~(212) already make the distribution at each $v_\parallel$ a passively mixed field. The normal fields of truncated drift-fluid and gyrofluid hierarchies \citep{Tassi2015,Tassi2017,Tassi2019,Tassi2023} sit at the zeros of a Hermite polynomial \citep{Tassi2022} and, except one field at $v_\parallel=0$ for an odd number of moments, remain coupled in three dimensions; \citet{ParkerDellar2015} identified the eigenvectors of the truncated Hermite system of linearised Vlasov--Poisson with discrete Case--Van Kampen modes, and \citet{Mandell2024} include the largest Gauss--Hermite node in the time step of the GX code. CMM showed that in gyrokinetics the distribution at $v_\parallel=0$ is advected by the gyroaveraged potential alone. S09 (\S6.3) showed that the kinetic equation is linear in the Lagrangian frame of the $\boldsymbol E\times\boldsymbol B$ flow, and \citet{Kunz2015} restated this for pressure-anisotropic plasmas. \citet{KanekarThesis2015} and the GANDALF paper \citep{Kanekar2025} describe the compressive distribution of KRMHD as a kinetic passive scalar, and \citet{Kanekar2014APS} studied a model kinetic passive scalar advected by a stochastic flow; none of these works decomposes the distribution into slices. Remark~\ref{sl:remlit} in section~\ref{sec:slices} gives further precedents, including two-dimensional ones. What is new here is the exact decomposition of the three-dimensional KRMHD compressive sector into passive slices streaming along the perturbed field lines, in the continuum and in truncated hierarchies, and its consequences.

A companion paper \citep{KanekarPrep} studies the Hermite spectrum of forced KRMHD turbulence. It rests on the frozen-line theorem, by which, when the induction equation is ideal, the volume-integrated single-point statistics of the compressive distribution evolve as in linear phase mixing along the initial field lines. That theorem is proved there without the slices; section~\ref{sec:discussion} gives its slice form.
